%% ===========================================================================
\FloatBarrier
\section{Numerical results and discussion}
\label{sec:results}
%% ===========================================================================

The numerical study examines window orientation and aspect ratio,
fractional order and relaxation time, specimen thickness, the resulting
parameter map and the response of the electrode configurations. The frequency, coupling and transient-response quantities are obtained from
the defined finite-element and reduced-order procedures. The preceding
benchmarks establish mesh-refinement behaviour; additional verification
tables are provided in the verification appendix.

\FloatBarrier
\subsection{Window orientation and window aspect ratio}
\label{sec:orientation}

Figure~\ref{fig:orientation} illustrates the influence of the orientation
of an anisotropic sampling window. Since the adopted gradient energy is linear
in the second-moment tensor $M$, rotation changes its components without
changing the strip geometry. The window is held at a
half-width of \SI{3}{\micro\meter} along its long axis; its aspect ratio is
${\rm AR}=a_1/a_2$. The strip dimensions are
\SI{200}{\micro\meter}$\times$\SI{10}{\micro\meter}, with poling along $x$; in the adopted local gradient model, the prescribed window determines
the directional second moments and therefore weights the components of the
strain-gradient stiffness. Over the $6\times19$ cuboid cases the extensional mode remains the one
with the largest coupling, its short-circuit frequency varying smoothly from
$4907$ to \SI{5296}{kHz}, consistent with tracking the same mode under window rotation.

\subsubsection{Symmetry of the orientation response}

The strip maps onto itself under $\theta\to\theta+180^\circ$. The second
invariance is a symmetry of the whole problem and not of the energy alone: the
gradient energy is \emph{linear} in $M$, so a change of sign of $M_{12}$ by
itself changes the response, and what leaves it unchanged is the reflection
$y\to-y$, which maps $\theta\to-\theta$ and at the same time maps the strip,
the electrode pair and -- for poling along $x$, the configuration plotted here
-- the material onto themselves. The response is therefore even in $\theta$ for
the poling-$x$ configuration; for poling along $y$ the same reflection reverses
the poling direction and is not a symmetry of the material, so only the period
$180^\circ$ holds and the orientation study is confined to poling $x$. Neither
invariance is assumed: both are evaluated at \SI{15}{\degree} intervals for the
extreme aspect ratios, where they reproduce the computed points to
$3\times10^{-11}$ relative.

\subsubsection{Dependence on the orientation angle}

The extensional frequency and coupling vary monotonically with $\theta$ and approach limiting values near
$\theta=90^\circ$. The extensional frequency rises from $5011$ to
\SI{5296}{kHz} at ${\rm AR}=2$ and from $4907$ to \SI{5287}{kHz} at
${\rm AR}=10$, that is by $6$--$8\,\%$, while the coupling falls from $0.4654$
to $0.4017$ and from $0.4874$ to $0.4032$ between the same two orientations, a
change of $14$--$17\,\%$. The fundamental mode of the same strip, which carries
a coupling three orders of magnitude smaller, moves from \SI{146.4}{kHz} to
\SI{158.1}{kHz} at ${\rm AR}=10$. At $\theta=0^\circ$, increasing the window aspect ratio from ${\rm AR}=2$
to ${\rm AR}=10$ changes the coupling from $0.4654$ to $0.4874$,
approximately $5\,\%$. At $\theta=90^\circ$, the computed values
converge to the narrow interval $0.4017$--$0.4032$, indicating weak
sensitivity to aspect ratio in the transverse orientation.

The plotted quantity is the total extensional-mode coupling, in which the
piezoelectric contribution is dominant; the flexoelectric contribution at
$\mu=10^{-6}$\,\si{\coulomb\per\meter} is of order $10^{-4}$,
so the window acts on the modal coupling through
the gradient stiffness rather than by adding a flexoelectric term directly.
Figure~\ref{fig:design} separates the two. Panel (a) is the change of the
coupling against the aligned window: with the largest change for the most elongated window,
$-17.3\,\%$ at ${\rm AR}=10$ and $\theta=90^\circ$, and the smallest change for the least elongated window, $-13.7\,\%$ at ${\rm AR}=2$. Panel (b) is
$(M_{11}-M_{22})/(M_{11}+M_{22})$, shown for ${\rm AR}=2$, which
passes through zero at $45^\circ$ and changes fastest there. The orientation dependence is governed by the second-moment tensor
within the adopted local model. Thus, changing the window orientation produces a $14$--$17\,\%$ variation
in the calculated coupling. When an elongated window is aligned transverse to
the gradient, further changes in its aspect ratio affect the coupling by no
more than $5\,\%$.

\begin{figure}[!htbp]
\centering
\includegraphics[width=\textwidth]{F4_orientation.pdf}
\caption{Variations of (a) the extensional frequency and (b) its coupling factor, both normalized to the aligned ${\rm AR}=2$ window, over the full circle of window orientation, and of (c) the coupling factor and (d) the frequency as maps over orientation and window aspect ratio, for a \SI{200}{\micro\meter}$\times$\SI{10}{\micro\meter} strip with poling along $x$, plane stress and a $24\times2$ mesh. Markers are computed points, lines a periodic spline. The long window half-width is \SI{3}{\micro\meter}.}
\label{fig:orientation}
\end{figure}

\begin{figure}[!htbp]
\centering
\includegraphics[width=0.86\textwidth]{F5_design.pdf}
\caption{Variations of (a) the change of the extensional coupling relative to the aligned window, one curve per window aspect ratio, and of (b) the anisotropy of the second-moment tensor at ${\rm AR}=2$, for the same length-poled strip and rotated windows. The percentage decrease relative to the aligned window ranges approximately from
$14\,\%$ to $17\,\%$ at $90^\circ$, a spread of about $3.6$ percentage points, and reaches $-17.3\,\%$ where the anisotropy of (b) is largest.}
\label{fig:design}
\end{figure}

\FloatBarrier
\subsection{Memory order, relaxation time and the damping map}
\label{sec:memorystudy}

Figure~\ref{fig:damping} shows the dependence of modal attenuation on
fractional order and the time-scale parameter in the reduced mechanical
stiffness-memory law. The fractional Kelvin--Voigt element of
\eqref{eq:kv} tends to a classical spring--dashpot model at $\alpha=1$;
orders below unity introduce a power-law memory response. The parameter
$\tau$ sets the dimensionless frequency $\omega\tau$ and therefore
scales the loss response. Unlike a standard linear solid, the Kelvin--Voigt
loss tangent does not exhibit a finite-frequency peak at $\alpha=1$. The convolution in \eqref{eq:step} uses either the
condensed open-circuit stiffness $K_{oc}$ or the mechanical-only
stiffness $K_{uu}$. These choices produce different damping responses.
The loss tangent \eqref{eq:tandelta} describes the frequency-domain
fractional modulus, whereas the plotted pole-based damping ratio comes
from the complex characteristic root; \eqref{eq:zetasmall} provides
only a restricted small-parameter approximation,
and the integer-order limit it is compared against is \eqref{eq:limitKV}.

\subsubsection{Influence of the memory stiffness operator}

Table~\ref{tab:memory} reports the ring-down of the extensional mode (mode 4,
\SI{6.855}{MHz}) at $\tau=\SI{6.9651}{\nano\second}$, that is $\omega\tau=0.3$,
with an initial tip displacement that produces \SI{143.78}{\volt} of
open-circuit voltage at $t=0$. The coupling values for this mode correspond to two stiffness
conventions:
$k^2_{4}(K_{uu})=0.4847$ is the mode-based measure on the elastic memory matrix
$K_{uu}$ and $k^2_{4}(K_{sc})=0.4784$ is the same measure on the short-circuit
stiffness $K_{sc}$, both on the $24\times2$ mesh of the ring-down; the third,
$0.4781$ from the application example is $k^2_{4}(K_{sc})$ again,
recomputed with $\mu=0$ because that example quotes the material without
flexoelectricity. The gap of $1.3\,\%$ between $0.4847$ and $0.4784$ reflects the
choice of short-circuit stiffness in the coupling measure, while the
$0.07\,\%$ between $0.4784$ and $0.4781$ reflects the inclusion of $\mu$.

The damping ratio is extracted from two successive peaks with
\eqref{eq:zeta}. At $\omega\tau=0.3$ the ratio of
two successive voltage peaks is $2.5947$, so $\delta=0.95345$ and
$\zeta_{\rm eff}=0.15003$ for $K_{oc}$ at $\alpha=1$, against the exact
$\omega\tau/2=0.15000$: the implementation reproduces the classical limit to
$0.02\,\%$. Refining the step gives $\zeta_{\rm eff}=0.150111$, $0.150057$ and
$0.150029$ at $N$, $2N$ and $4N$ steps over twenty periods, so the residual is
first order in $\Delta t$ (observed order $0.97$) and extrapolates to
$0.150000$; every ring-down quoted in this section is run at the $4N$ step.
Replacing $K_{oc}$ by $K_{uu}$ \emph{in the memory operator only}
removes the electromechanical stiffening from the remembered term and gives
$\zeta(K_{uu})/\zeta(K_{oc})=0.431$ at $\alpha=1$, a $56.9\,\%$
reduction in the reported damping ratio for this configuration. This
comparison does not represent a short-circuited electrode problem. The
difference is of the order of the coupling ($k^2_{4}=0.4847$) but does
\emph{not} reduce to the approximation $1-k^2=0.5153$: the two memories differ by
the piezoelectric stiffness that the condensation puts into $K_{oc}$, and
its effect on the ring-down has to be computed rather than estimated. Thus the stiffness matrix used in the memory law must be specified
when reporting damping.

Neither simulated memory operator is derived from dielectric relaxation.
To distinguish the physics, suppose instead that a common scalar dielectric
relaxation factor $g(s)=1+(s\tau)^\alpha$ multiplies the dielectric block,
while the piezoelectric coupling is instantaneous. With identical electrode
constraints and a fixed reference potential, Schur elimination gives
$C_\phi(s)=g(s)C_\phi^0$ and
$K_{\rm cond}(s)=K_{uu}+g(s)^{-1}B(C_\phi^0)^{-1}B^{\mathsf T}$.
The \emph{inverse} factor is required because the dielectric block is
inverted; dielectric relaxation is therefore neither of the two mechanical
stiffness-memory models evaluated here. The measured ratio $0.431$ is
specific to the proportional Kelvin--Voigt test at $\alpha=1$ and to the
chosen specimen and electrodes.

%% The memory-matrix comparison, moved from the appendix to the text.
\begin{table}[!htbp]
\centering
\small
\caption{Ring-down with the memory carried on the condensed stiffness
$K_{\rm eff}$ and on the mechanical-only stiffness $K_{uu}$ (mode 4,
\SI{6.855}{MHz}, $\omega\tau=0.3$). At $\alpha=1$ the entry is the damping
ratio \eqref{eq:zeta}; at $\alpha=0.7$ the decay is not exponential and the
entry is the first-period equivalent damping index derived from the
logarithmic decrement, rather than the decrement itself.}
\label{tab:memory}
\begin{tabular}{@{}lcc@{}}
\toprule
memory matrix & $\alpha=1$: $\zeta_{\rm eff}$ &
  $\alpha=0.7$: first-period $\zeta_{\rm eq}$ \\
\midrule
$K_{\rm eff}$ & $0.15003$ & $0.1942$ \\
$K_{uu}$      & $0.06465$ & $0.0796$ \\
\midrule
ratio $K_{uu}/K_{\rm eff}$ & $0.431$ & $0.410$ \\
\bottomrule
\end{tabular}
\end{table}


\subsubsection{Decay at fractional order}

At $\alpha=0.7$ the free decay is not exponential and a single damping ratio
does not exist: the successive decrements of the $K_{oc}$ ring-down are
$1.244$, $1.547$, $0.010$, $0.403$, $0.240$ and $0.151$; the first, second,
fourth, fifth and sixth of them are unchanged to three digits between
$\Delta t$, $\Delta t/2$ and $\Delta t/4$, so they are a property of the
nonexponential fractional response rather than the time step. The third
decrement reflects a near-cancellation
between two peaks of comparable height. A straight line fitted to $\ln A_k$, with $A_k$ the successive peak amplitudes,
over the twenty periods therefore returns a number that is not a damping ratio.
Accordingly, the $\alpha=0.7$ result is reported as a first-period
decrement rather than as a fitted exponential slope. The $\alpha=1$ entries are
unaffected: there the twenty successive decrements agree to $7\times10^{-10}$
relative, which is why \eqref{eq:zeta} applies.

\subsubsection{Kelvin--Voigt/Zener model-form sensitivity}

This comparison uses the modulus magnitudes of two model families:
$C^*_{KV}/C_0=1+z$ and $C^*_{Z}/C_0=(1+r z)/(1+z)$,
where $z=(\mathrm{i}\omega\tau)^\alpha$ and $r=C_\infty/C_0=1.1$.
The same $C_0$ and $\tau$ are prescribed, not fitted to matched experimental spectra. At the device point $\omega\tau=0.3$ the
Kelvin--Voigt modulus magnitude exceeds the Zener magnitude by $3.5\,\%$ at $\alpha=1$ and by
$22.6\,\%$ at $\alpha=0.7$; at $\omega\tau=1$ the ratio of the two is $1.35$
and $1.62$, respectively. At $\omega\tau=100$ and $\alpha=1$ the ratio is approximately $91$. Thus, a single-relaxation fit calibrated over a narrow frequency interval
cannot be assumed accurate across the full range of the design map.

\subsubsection{Fractional-order damping map}

Panel~(a) traces a pole of the proportional modal displacement transfer
function, whose principal-branch characteristic equation is
$s^2+\omega_0^2[1+(s\tau)^\alpha]=0$, in the complex plane, traced over $\omega\tau$ for seven relaxation times and colored by
$\tau$, with markers at five Caputo orders; panel~(b) gives the
damping ratio $\zeta=-\mathrm{Re}\,s/|s|$, where $s$ is the dimensional
complex pole and $\omega_0$ the undamped angular frequency, over the
$(\alpha,\omega\tau)$ plane. Both are computed with branch continuation: the
root is followed from one $\omega\tau$ to the next by initializing the Newton
iteration at an adjacent converged pole. The characteristic residual is checked;
unconverged Newton iterates are rejected. Fractional responses also contain
a branch-cut contribution, so pole damping is not a global exponential decay rate.
For nonproportional memory the characteristic equation is
$\det[s^2M+K_{oc}+\tau^\alpha s^\alpha K_d]=0$. All computed poles lie in the left half-plane
for the sampled $(\alpha,\tau)$ values, consistent with decay for
these cases; this finite parameter grid is not a general stability proof. The integer-order case sits at
$|\mathrm{Im}\,s|/\omega_0\simeq1$ for $\omega\tau\lesssim0.1$ and moves down
and to the right as $\omega\tau$ grows, so the dashpot lowers the frequency and
eventually over-damps the mode: at $\omega\tau=10$ the imaginary part has
vanished and $\zeta=1$. The damped frequency need not increase monotonically with $\tau$ for
$\alpha<1$. For example, at $\alpha=0.9$, independent characteristic-root
calculations give $\operatorname{Im}(s)/\omega_0=1.00223,1.01571,1.01207$
and $0.98457$ at $\omega_0\tau=0.02,0.30,0.60$ and $1.00$, respectively.
The frequency first rises and then falls within this interval. At
$\omega\tau=10$ and $\alpha=0.5$ the damped frequency is $2.13\,\omega_0$
($2.45\,\omega_0$ at $\alpha=0.7$), against $\mathrm{Im}\,s=0$ for $\alpha=1$ at
the same $\omega\tau$. The damping ratio is not generally monotone in fractional order. At
fixed $\omega_0\tau=0.3$, the independently checked pole ratios for
$\alpha=0.3,0.5,0.7,0.9,1$ are respectively
$0.1077,0.1586,0.1794,0.1661,0.1500$, with an intermediate maximum.
Dependence on $\omega_0\tau$ must likewise be assessed at fixed $\alpha$;
in the integer-order case sufficiently large $\omega_0\tau$ leads to
over-damping, whereas fractional orders exhibit different pole trajectories. Such broad relaxation spectra also arise in dielectric loss models,
but the fractional operator simulated here is applied to mechanical stiffness,
not to a frequency-dependent dielectric permittivity. The band $\omega\tau\le1$ corresponds to
$\tau\lesssim\SI{23.6}{\nano\second}$ at \SI{6.73}{MHz}.
The largest plotted value $\tau=\SI{32}{\nano\second}$ instead
corresponds to $\omega\tau\simeq1.35$ and extends beyond that
restricted band. These are numerical parameter studies rather than
experimentally calibrated damping predictions. Both $\alpha$ and
$\tau$ influence the resulting pole and decay.

\begin{figure}[!htbp]
\centering
\includegraphics[width=\textwidth]{F6_damping.pdf}
\caption{Variations of (a) the pole of a mode at $f_0=\SI{6.73}{MHz}$ ($\omega_0=2\pi f_0$) over $\omega\tau=0.02$--$1.35$ for seven relaxation times, identified by $\tau$ in the legend and distinguished by line style and marker, with markers at five Caputo orders, and of (b) the damping ratio over $\omega\tau=10^{-1.5}$--$10^{1.5}$ and $\alpha=0.2$--$1$ with labeled contours, for the fractional Kelvin--Voigt element with branch continuation. The dashed arc is the exact Kelvin--Voigt locus, on which every $\alpha=1$ pole lies. Fractional trajectories depend on order; increasing relaxation time does
not universally increase their damped frequency. Integer-order poles
become overdamped for sufficiently large $\omega_0\tau$.}
\label{fig:damping}
\end{figure}

\FloatBarrier
\subsection{Window size and shape against the second moment}
\label{sec:windowshape}

Figure~\ref{fig:window} compares window size and geometry for a
benchmark with a known classical solution. Because the adopted gradient
energy depends on $M$, windows with identical second moments give
identical predictions within this local model. Panel~(a) evaluates this relationship on an isotropic plane-strain test problem (rectangle
$1\times0.2$, $20\times4$ mesh, uniaxial edge traction, displacement averaged
over the loaded edge), plotted against the normalized window size for an
ellipsoidal window of semi-axis $l$, a cuboidal window of half-width $l$, and a
cuboidal window scaled to carry the same $M$ as the ellipsoid,
$a=l\sqrt{3/5}$. The third curve is indistinguishable from the first: within
the adopted second-moment model the shape of the window is irrelevant at fixed
$M$. The second curve shows the difference caused by using equal half-widths rather than equal moments: equal
half-widths instead of equal $M$ over-stiffen the response, and the
discrepancy grows with the window size. Writing $L_x$ for the rectangle length, at $l/L_x=0.2$ a cuboid of half-width
$l$ retains $0.922$ of the classical loaded-edge displacement, an ellipsoid of
semi-axis $l$ retains $0.941$, and the equal-$M$ cuboid agrees with the ellipsoid
curve to approximately three decimal places. At fixed $M$, the reported displacements for
equal-second-moment windows agree within $10^{-3}$, which is consistent with the $5/3$ ratio of the corresponding second-moment
coefficients. The assumed operating range is
$a/L\le0.2$, the largest value plotted.

Panels (b) and (c) present the corresponding static displacement fields. The baseline clamped strip is solved under an end load for three window sizes and
each solution is compared with the classical one of the same geometry and load.
The gradient terms stiffen the member. The difference against the classical
solution vanishes at the clamped edge, where both fields satisfy the same
kinematic conditions, and then accumulates along the member instead of decaying
over a thin layer: at these window sizes the gradient effect is distributed
rather than local, and its size grows with the window. That is why $a/L\le0.2$
is quoted as an assumed operating range and not as a validated truncation
threshold. Thus, at a fixed second moment, the plotted responses agree to the reported
precision; window size and anisotropy remain the parameters represented in the
adopted second-moment model. The comparison uses the ellipsoidal second moment \eqref{eq:Mellip}
and its rotation \eqref{eq:Mrot}; the classical profile follows from
\eqref{eq:limitM} with flexoelectricity suppressed.

\begin{figure}[!htbp]
\centering
\includegraphics[width=\textwidth]{F7_window.pdf}
\caption{Variations of (a) the loaded-edge displacement of an isotropic plane-strain rectangle against normalized window size, for an ellipsoid, a cuboid of equal half-width and a cuboid of equal second moment, and of (b) the deflected shape of a clamped strip for three window sizes together with (c) its difference from the classical profile. Windows of equal second moment collapse onto one curve. Panel (a) uses a $1\times0.2$ rectangle and a $20\times4$ mesh under uniaxial edge traction; panels (b,c) use the baseline clamped strip with $a/L=0.05$, $0.1$ and $0.2$.}
\label{fig:window}
\end{figure}

\FloatBarrier
\subsection{Thickness: a flexoelectric size effect against a thickness-independent
piezoelectric coupling}
\label{sec:thickness}

Figure~\ref{fig:thickness} examines the thickness dependence of the
piezoelectric and flexoelectric coupling factors: for a prescribed pure-bending curvature $\kappa_c$, the transverse strain
gradient has order $\kappa_c$, independent of $h$. In this electrode and
modal setting the normalized flexoelectric coupling can approach an $h^{-2}$
trend when the window is much thinner than the strip, because mechanical and
electrical normalization also change with $h$; the gradient itself does not
scale as $\kappa_c/h$. The strip is analyzed at eight
thicknesses between \SI{10}{\micro\meter} and \SI{1}{\micro\meter} with a
\SI{200}{\micro\meter} length and the window held at half-widths
$3\times1$\,\si{\micro\meter}, so that the thickness is the only variable. The complete sweep is tabulated
in Table~\ref{tab:thickness}, with supporting mesh data in \ref{app:verification}. The thin configurations require a qualification: the
window half-thickness is \SI{1}{\micro\meter} throughout, so at
$h=\SI{2}{\micro\meter}$ it is half the strip and at $h=\SI{1}{\micro\meter}$
it equals the strip. The second-moment tensor of the full window is used at
every thickness and no truncated-window boundary correction is applied, so $M$
acts there as a local effective gradient tensor rather than as the moment of a
window small compared with the body. The two panels are two electrical
configurations of that strip: panel~(a) is poled through the thickness
(electrodes on the long faces, gap $h$, $\kappa_{yy}=\kappa_{33}$) and
panel~(b) along the length (electrodes on the short faces, gap $L$,
$\kappa_{xx}=\kappa_{33}$).

\begin{figure}[!htbp]
\centering
\includegraphics[width=\textwidth]{F8_thickness.pdf}
\caption{Variations of (a) the flexoelectric bending coupling against strip thickness at eight thicknesses, with local log--log slopes, a $96\times4$ mesh comparison at both ends and the $h^{-2}$ guide, of (b) the piezoelectric extensional coupling of the same strip with $\mu=0$, constant to $0.07\,\%$ over $10\to1\,\si{\micro\meter}$, and of (c) their ratio, which spans $2.9\times10^{-4}$ to $7.0\times10^{-3}$.}
\label{fig:thickness}
\end{figure}

\subsubsection{The flexoelectric size effect}

The flexoelectric coupling rises by a factor $24$ from $1.67\times10^{-4}$ at
\SI{10}{\micro\meter} to $4.0\times10^{-3}$ at \SI{1}{\micro\meter}. Its local
log--log slope is $-1.96$ at the thick end, close to the $h^{-2}$ scaling of a
gradient effect with a window much smaller than the thickness, and falls to
$-1.53$, $-1.27$, $-0.88$ and $-0.44$ on the successively thinner pairs, the seven
segments spanning $-1.96$ to $-0.44$. The reduced slope persists under mesh refinement: a
$96\times4$ mesh reproduces the $48\times4$ values to $0.024\,\%$ and
$0.38\,\%$ at the two ends of the range, and the residual at $\mu=0$ is
$1.0\times10^{-10}$, seven orders of magnitude below the signal. A dimensionless group associated with the observed saturation is $a_2/h$, the window
half-thickness over the strip thickness, which runs from $0.10$ at
$h=\SI{10}{\micro\meter}$ to $1.0$ at \SI{1}{\micro\meter}: the slope decreases
monotonically as $a_2/h$ grows in this numerical model. At the thin end,
however, the nominal untruncated window is no longer contained inside the
specimen: the second-moment tensor must be interpreted as an effective local
gradient parameter rather than a literal average over the surrounding body. The
same stiffening appears in the fundamental frequency, which follows
$210.4,108.0,50.2,34.7$\,kHz where pure Euler--Bernoulli scaling from the
thickest strip would give $210,105,42,21$\,kHz.

Figure~\ref{fig:modetrack} reports modal tracking across the thickness
sweep, supporting a consistent sampled mode branch. Shapes are pulled back
by $y=(h/h_{\rm ref})y_{\rm ref}$ before comparison with $h_{\rm ref}=\SI{10}{\micro\meter}$:
the physical $y$ and $xy$ Hermite derivative coordinates are multiplied by
$h/h_{\rm ref}$, while values and $x$ derivatives are unchanged. Using the
reference mass matrix, the minimum modal assurance criterion is $0.999208$
and the largest off-diagonal value is $2.47\times10^{-6}$. The ratio $f_2/f_1$
lies between $6.21$ and $6.27$ at the four tested thicknesses. These samples
do not prove absence of crossings between them or certify the small coupling shifts.

\begin{figure}[!htbp]
\centering
\includegraphics[width=0.90\textwidth]{FB1_tracking.pdf}
\caption{Variations of (a) the defect $1-\mathrm{MAC}$, the MAC being the modal assurance criterion, between the mode shape at each thickness and the shape at $h=\SI{10}{\micro\meter}$, with the MAC printed, and of (b) the frequencies of the first two modes and their ratio, over a sweep from $10$ to $1\,\si{\micro\meter}$. Shapes use the normalized-thickness pullback and reference mass matrix; the sampled MAC defect is below $10^{-3}$. The same strip, electrode configuration and thickness values are used for the mode-tracking analysis.}
\label{fig:modetrack}
\end{figure}

Figure~\ref{fig:modes} shows the displacement and potential fields
of the tracked bending mode. The surface is the computed
open-circuit bending displacement, colored by the solved electric potential,
and the inset shows the same potential contours on the strip coordinates.

\begin{figure}[!htbp]
\centering
\includegraphics[width=\textwidth]{FB2_modes.pdf}
\caption{Open-circuit bending mode and its potential field. The three-dimensional surface is the computed transverse displacement in nanometers, colored by the solved electric potential; the inset shows labeled potential contours on the strip coordinates, with contour labels in millivolts. Both views use one finite-element solution for the baseline open-circuit bending configuration.}
\label{fig:modes}
\end{figure}

%% generated table -- do not edit (see AUTHOR_TODO.md)
\begin{table}[!htbp]
\centering
\small
\caption{Thickness sweep behind Fig.~\ref{fig:thickness}. First frequency,
flexoelectric bending coupling and extensional coupling of the poled strip. The
bending column carries $\mu=10^{-6}$\,\si{\coulomb\per\meter}, the extensional
column $\mu=0$. These are the archived sweep results. Convergence of the extensional
coupling does not establish convergence of the first bending frequency
or of the small flexoelectric spectral shift at $h=\SI{1}{\micro\meter}$.}
\label{tab:thickness}
\begin{tabular}{@{}r r r r@{}}
\toprule
$h$ (\si{\micro\meter}) & $f_1$ (kHz) & $k^2_{\rm flexo}$ & $k^2_{\rm piezo}$ \\
\midrule
10 & 209.0 & 1.668e-04 & 0.57268 \\
7 & 147.7 & 3.358e-04 & 0.57283 \\
5 & 107.3 & 6.307e-04 & 0.57294 \\
3.5 & 77.7 & 1.167e-03 & 0.57301 \\
2.7 & 62.5 & 1.735e-03 & 0.57305 \\
2 & 49.9 & 2.541e-03 & 0.57307 \\
1.4 & 40.2 & 3.481e-03 & 0.57309 \\
1 & 34.6 & 4.040e-03 & 0.57310 \\
\bottomrule
\end{tabular}
\end{table}


\subsubsection{The piezoelectric coupling is thickness independent}

For the extensional configuration, the
quasi-static coupling takes the values $0.57268$, $0.57294$,
$0.57307$ and $0.57310$ at $10$, $5$, $2$ and \SI{1}{\micro\meter}, that is it
changes by $0.07\,\%$ between \SI{10}{\micro\meter} and \SI{1}{\micro\meter},
against the constitutive closed form $0.5752$ for the same material and the
same plane-stress kinematics. The computed extensional piezoelectric coupling remains nearly unchanged
across this thickness interval, whereas the flexoelectric bending coupling
increases by a factor of $24$. In this model, the predicted relative
flexoelectric contribution therefore becomes larger for thinner strips; the
crossover thickness depends on the window and on $\mu$, rather than on the piezoelectric
constants.

\subsubsection{Piezoelectric contribution to bending response}

The bending configuration is poled along $y$ with electrodes on the long faces,
and the crystal carries $e_{31}$ and $e_{33}$ as well, allowing the two contributions to be assessed separately. With $\mu=0$ the piezoelectric coupling of that
mode is $1.0\times10^{-10}$ at every thickness, because the piezoelectric charge
$e_{31}\varepsilon_{xx}+e_{33}\varepsilon_{yy}$ is odd in $y-h/2$ and its
integral over a face electrode vanishes, so the arrangement carries no net
piezoelectric charge. With $e=0$ the flexoelectric coupling runs from
$1.76\times10^{-4}$ to $5.5\times10^{-3}$, a factor $31$, while the full crystal
gives the $24$ reported above: the piezoelectric terms shift the flexoelectric
coupling by up to $26\,\%$ at the thinnest strip, through the elastic stiffness
of the same mode rather than through a charge they contribute. The reported headline value refers to the full-crystal calculation.

\subsubsection{Numerical precision at the thinnest geometry}

The $h=\SI{1}{\micro\meter}$ point exhibits the greatest conditioning
sensitivity in the reported sweep. The flexoelectric coupling is inferred
from a small electrode-induced modal shift in an ill-conditioned generalized
eigenvalue problem. The corrected-model sensitivity record comprises
\emph{one source-tagged numerical environment} (Python 3.13.5, NumPy
2.3.5, SciPy 1.17.0) evaluated using four alternative eigenvalue
algorithms, rather than four independently verified software environments.
The four paths give $k^2=4.0342\times10^{-3}$--$4.0400\times10^{-3}$,
a range-over-midpoint spread of $0.143\,\%$ at $h=\SI{1}{\micro\meter}$;
at $h=\SI{10}{\micro\meter}$ the spread is
$3.4\times10^{-6}\,\%$. Historical numerical-stack records computed
with the unreduced plane-stress material are excluded from this comparison;
the values do not constitute independent cross-platform replication.
 Two significant digits are used for presentation, but are not a certified
error bound: fresh finer-mesh checks still exhibit conditioning-sensitive
eigenpair residuals (Sec.~\ref{sec:limitations}). Hence
thinning the strip multiplies the flexoelectric bending coupling by $24$ while
leaving the piezoelectric extensional coupling unchanged to $0.07\,\%$, showing the different thickness sensitivities of the two coupling
mechanisms in the adopted model. The estimate \eqref{eq:xflexo} is the $h^{-2}$ guide drawn in
panel~(a), and the inverse relation \eqref{eq:muV} provides a possible route for estimating $\mu$ from measured electrical response.

\FloatBarrier
\subsection{The design map}
\label{sec:designmap}

Figure~\ref{fig:designmap} compares the flexoelectric bending and
piezoelectric extensional couplings across the same thickness and window
parameter grid. The former depends strongly on both parameters, whereas
the latter changes only weakly over the investigated range. Panel~(a) is the flexoelectric
bending coupling on a grid of eight thicknesses and eight window
half-thicknesses, computed with the mode-based measure \eqref{eq:k2mode} and
the same $48\times4$ mesh as the thickness sweep; panel~(b) is the
ratio of that bending coupling to the mode-free extensional coupling
\eqref{eq:kstatic} of the same strip. The dotted line $2a_2=h$ separates
windows whose full transverse extent can fit within the strip from those
that extend beyond its faces; the hatched region represents extrapolation
of the prescribed local second-moment model, not a boundary-aware average.
The computed grid extends to $a_2=\SI{32}{\micro\meter}$, giving
$a_2/L=0.16$ at its top boundary, within the prescribed
range $a_2/L\le0.2$.

For geometric interpretation, the $8\times8$ grid must be partitioned:
25 of its 64 points satisfy $2a_2\le h$, a necessary condition for the
full centered window to fit at the strip mid-plane; 39 points do not.
For those 25 mid-plane-fitting points, the computed flexoelectric bending
coupling spans $1.28\times10^{-4}$ to $1.01\times10^{-2}$ and the
flexoelectric-to-piezoelectric coupling ratio spans
$2.23\times10^{-4}$ to $1.75\times10^{-2}$. The corresponding ranges
for the 39 extrapolated points are $3.58\times10^{-6}$ to
$4.04\times10^{-3}$ and $6.25\times10^{-6}$ to $7.05\times10^{-3}$.
This necessary mid-plane condition does not validate a literal untruncated
averaging kernel at all positions: a window centered sufficiently near a
free face still requires a boundary-aware prescription even when $2a_2\le h$.
All 64 solutions are generated by the same prescribed local-gradient
constitutive model, not by a finite-domain integral kernel. These
partitions and values are independently reproducible with
the archived grid values and the accompanying geometric audit.

The design map shows three principal trends. First, the coupling is controlled by
the ratio $a_2/h$ rather than by the thickness alone. At
$h=\SI{10}{\micro\meter}$ shrinking the window half-thickness from $1$ to
$0.25\,\mu$m changes the coupling from $1.7\times10^{-4}$ to
$1.6\times10^{-4}$, a change of $5.4\,\%$, because the window is far thinner
than the strip in both cases; at $h=\SI{1}{\micro\meter}$ the same change of
window takes the coupling from $4.0\times10^{-3}$ to $1.0\times10^{-2}$, a
factor $2.5$. Second, the window dimension has the larger influence when it becomes
comparable to the strip thickness: over the complete local-model map (including extrapolated cases) the coupling spans
$3.6\times10^{-6}$ to $1.0\times10^{-2}$, of which the thickness alone
accounts for the previously reported factor of $24$. Third, the two
couplings never become comparable at these parameters: the ratio in panel~(b)
spans $6.3\times10^{-6}$ to $1.8\times10^{-2}$, and it reaches $0.1\,\%$ only
below $h=5.3\,\mu$m for the window of the paper, $a_2=\SI{1}{\micro\meter}$.
The map values are reported to two significant digits because the
thin-strip modal calculations remain sensitive to spectral conditioning;
no claim of design admissibility or specimen-level calibration is inferred
from the mid-plane-fitting geometry subset alone. The underlying
couplings are evaluated through a rank-one secular shift rather than direct
subtraction of open- and short-circuit eigenvalues.
Within the investigated parameter range, window size controls the
assigned gradient stiffness and strip thickness controls the size-dependent
contribution. Extrapolation to stronger flexoelectric coupling requires
additional material-specific assessment. The map distinguishes
the effect of thickness from the suppression caused by the gradient stiffness.

\begin{figure}[!htbp]
\centering
\includegraphics[width=\textwidth]{F11_designmap.pdf}
\caption{Variations of (a) the flexoelectric bending coupling and of (b) its ratio to the piezoelectric extensional coupling over strip thickness and window half-thickness, for $L=\SI{200}{\micro\meter}$ and $a_1=\SI{3}{\micro\meter}$, with the $0.01\,\%$, $0.1\,\%$ and $1\,\%$ contours. The dotted line marks $2a_2=h$; hatching identifies parameter combinations where a full centered window extends beyond the strip faces and the local-gradient interpretation is phenomenological. The grid ends at $a_2/L=0.16$, within the assumed but unvalidated $a_2/L\le0.2$ range.}
\label{fig:designmap}
\end{figure}

\FloatBarrier
\subsection{Transient response and the memory model}
\label{sec:transient}

Figure~\ref{fig:ringdown} presents the time-domain free-decay response of
the extensional mode, complementing the frequency-domain damping analysis. The strip starts from a tip displacement of
\SI{10}{\nano\meter} and is followed for eight Caputo orders between
$\alpha=1$ and $0.4$ at fixed $\tau=\SI{3}{\nano\second}$
($\omega\tau=0.13$): the tip displacement in panel~(a) and the open-circuit
voltage in panel~(b). The initial voltage, \SI{144}{\volt}, is set by the mode
shape and the electrode configuration and is independent of $\alpha$, as it
must be, since the memory term vanishes at $t=0^{+}$; the decay rate, by
contrast, depends on $\alpha$. The $\alpha=1$ curve follows the exponential decay of the classical
Kelvin--Voigt element, whereas the fractional-order responses are
nonexponential. Across the eight orders plotted, the
decrement of the tip displacement over the first two peaks, converted with
\eqref{eq:zeta}, increases from $0.065$ at $\alpha=1$ to $0.105$ at
$\alpha=0.6$ before decreasing to $0.085$ at $\alpha=0.4$ for the eight fractional orders examined.
For $\alpha<1$ this is a first-period equivalent damping index, not a single
exponential damping ratio, because the calculated fractional decay is nonexponential.
Using the integer-order early decay as representative underestimates the
first-period equivalent damping index by $62\,\%$ relative to its $\alpha=1$
value at $\alpha=0.6$. This trend applies to the tested mode and relaxation time; other modes
or parameter ranges may behave differently. A single scalar decrement generally underdetermines $(\alpha,\tau)$.
A full ring-down history can in principle identify both under a specified
model with known other parameters, but finite noisy data require an
identifiability and uncertainty study. The march is the one of \eqref{eq:newmark} with the
L1 history, and the decrement is converted with \eqref{eq:zeta}.

\begin{figure}[!htbp]
\centering
\includegraphics[width=\textwidth]{F9_ringdown.pdf}
\caption{Variations of (a) the tip displacement and (b) the open-circuit voltage of the extensional mode at \SI{6.855}{MHz} for eight Caputo orders, with modal memory on $K_{oc}$, eight modes, $\tau=\SI{3}{\nano\second}$ ($\omega\tau=0.13$), an initial tip displacement of \SI{10}{\nano\meter} and $4000$ L1 steps over ten periods. Insets in both panels resolve the first \SI{0.35}{\micro\second} using the same histories and axis units as the main panels. Only the $\alpha=1$ decay is exponential.}
\label{fig:ringdown}
\end{figure}

\FloatBarrier
\subsection{Application: a \SI{200}{\micro\meter} by \SI{10}{\micro\meter}
strip in open-circuit sensing}
\label{sec:application}

The worked case is a \SI{200}{\micro\meter} long, \SI{10}{\micro\meter} thick
$\mathrm{BaTiO_3}$ strip clamped on one short edge, poled along its length and
electroded on the two ends: the strain in the active region is extensional, so
the piezoelectric coupling dominates and the flexoelectric term is a
size-dependent correction. Table~\ref{tab:application} collects the four
polarization/kinematics combinations of that strip, each solved from the
condensed system \eqref{eq:saddle} with the electrode charge
\eqref{eq:electrode}, and each quoted using the mode-based coupling \eqref{eq:k2mode}. The lowest-frequency mode of the
extensional configuration is at \SI{147.73}{kHz}, above the clamped-free value
of a beam of the same length because the gradient terms stiffen the member,
while the strongly coupled extensional mode sits at \SI{6853.6}{kHz}.

\begin{table}[!htbp]
\centering
\small
\caption{The coupling of the strongest mode, its frequency and the
open-circuit voltage at a maximum nodal displacement magnitude of \SI{10}{\nano\meter}, for a
\SI{200}{\micro\meter}$\times$\SI{10}{\micro\meter} strip in four polarization
and kinematics combinations. The two plane-stress rows correspond to lengthwise and transverse poling,
and the couplings of the four rows are not additive. All selected modes are predominantly longitudinal (abbreviated longit.); their mass-weighted longitudinal displacement shares exceed $99.98\,\%$.}
\label{tab:application}
% table generated from the archived results
\begin{tabular}{@{}llcccc@{}}
\toprule
poling & kinematics & mode & $f$ (kHz) & $k^2$ & $V$ at \SI{10}{\nano\meter} (V) \\
\midrule
y & plane strain & longit. (4) & 6721.8 & 0.1749 & 3.48 \\
y & plane stress & longit. (4) & 6323.8 & 0.1036 & 2.22 \\
x & plane strain & longit. (5) & 7161.0 & 0.4024 & 141.45 \\
x & plane stress & longit. (5) & 6853.6 & 0.4781 & 143.71 \\
\bottomrule
\end{tabular}

\end{table}

The results show the following differences between configurations.

\begin{enumerate}
\item Poling along the length
  gives $k^2=0.4781$ and \SI{143.71}{\volt} at \SI{10}{\nano\meter} of motion;
  transverse poling with face electrodes gives $0.1036$ and \SI{2.22}{\volt},
  approximately $65$ times less voltage and $4.6$ times less $k^2$.
  Both selected plane-stress modes are predominantly longitudinal; this is
  a comparison of poling, reduced tensors and electrode arrangements,
  not a comparison of longitudinal and bending mechanical modes.
\item The quasi-static coupling is $0.5727$ and thickness independent:
  the relevant measure below the first resonance and over the
  investigated thickness range.
\item The damping is set by $(\alpha,\tau)$ and by the memory matrix,
  not by a material constant: at $\omega\tau=0.3$ and $\alpha=1$,
  $\zeta_{\rm eff}=0.15003$ on $K_{oc}$ against $0.0646$ on $K_{uu}$, a
  ratio of $0.431$. A damping number quoted without $\alpha$, $\tau$ and the
  memory matrix is under-specified.
\end{enumerate}

The results depend on prescribed parameters; calibration routes and
applicability limits are specified in Sec.~\ref{sec:calibration}.

\FloatBarrier
